\documentclass[10pt,a4paper]{amsart}
\usepackage{amsmath,amssymb,amsthm,graphicx}
% Retained preamble for the Stacks Project chapter snapshot in this directory.
% The chapter's own preamble.tex is a copy of the shared Stacks Project preamble; the block
% below is sliced MECHANICALLY, line for line, from it. Nothing was added, removed or
% redefined; the wrapper class, the xr-hyper/hyperref wiring and the externaldocument list
% are omitted because the assembler supplies the document class.
%
% lines taken from preamble.tex: 3, 5, 9, 11, 12, 156, 157, 158, ... (51 lines total)
%
% package coverage restored from the chapter's own preamble (2026-10-13): xy, xy-2cell

\usepackage[all]{xy}
\xyoption{2cell}
\UseAllTwocells

\newenvironment{reference}{\comment}{\endcomment}
\newenvironment{slogan}{\comment}{\endcomment}
\newenvironment{history}{\comment}{\endcomment}
\theoremstyle{plain}
\newtheorem{theorem}[subsection]{Theorem}
\newtheorem{proposition}[subsection]{Proposition}
\newtheorem{lemma}[subsection]{Lemma}
\theoremstyle{definition}
\newtheorem{definition}[subsection]{Definition}
\newtheorem{example}[subsection]{Example}
\newtheorem{exercise}[subsection]{Exercise}
\newtheorem{situation}[subsection]{Situation}
\theoremstyle{remark}
\newtheorem{remark}[subsection]{Remark}
\newtheorem{remarks}[subsection]{Remarks}
\numberwithin{equation}{subsection}
\def\lim{\mathop{\mathrm{lim}}\nolimits}
\def\colim{\mathop{\mathrm{colim}}\nolimits}
\def\Spec{\mathop{\mathrm{Spec}}}
\def\Hom{\mathop{\mathrm{Hom}}\nolimits}
\def\Ext{\mathop{\mathrm{Ext}}\nolimits}
\def\SheafHom{\mathop{\mathcal{H}\!\mathit{om}}\nolimits}
\def\SheafExt{\mathop{\mathcal{E}\!\mathit{xt}}\nolimits}
\def\Sch{\mathit{Sch}}
\def\Mor{\mathop{\mathrm{Mor}}\nolimits}
\def\Ob{\mathop{\mathrm{Ob}}\nolimits}
\def\Sh{\mathop{\mathit{Sh}}\nolimits}
\def\NL{\mathop{N\!L}\nolimits}
\def\CH{\mathop{\mathrm{CH}}\nolimits}
\def\proetale{{pro\text{-}\acute{e}tale}}
\def\etale{{\acute{e}tale}}
\def\QCoh{\mathit{QCoh}}
\def\Ker{\mathop{\mathrm{Ker}}}
\def\Im{\mathop{\mathrm{Im}}}
\def\Coker{\mathop{\mathrm{Coker}}}
\def\Coim{\mathop{\mathrm{Coim}}}
\DeclareMathSymbol{\boxtimes}{\mathbin}{AMSa}{"02}
\def\QCohstack{\mathcal{QC}\!\mathit{oh}}
\def\Cohstack{\mathcal{C}\!\mathit{oh}}
\def\Spacesstack{\mathcal{S}\!\mathit{paces}}
\def\Quotfunctor{\mathrm{Quot}}
\def\Hilbfunctor{\mathrm{Hilb}}
\def\Curvesstack{\mathcal{C}\!\mathit{urves}}
\def\Polarizedstack{\mathcal{P}\!\mathit{olarized}}
\def\Complexesstack{\mathcal{C}\!\mathit{omplexes}}
\def\Pic{\mathop{\mathrm{Pic}}\nolimits}
\def\Picardstack{\mathcal{P}\!\mathit{ic}}
\def\Picardfunctor{\mathrm{Pic}}
\def\Deformationcategory{\mathcal{D}\!\mathit{ef}}
\usepackage{cleveref}
\crefname{originalequation}{original equation}{original equations}
\Crefname{originalequation}{Original equation}{Original equations}
\makeatletter
\newlabel{equation-three-maps}{{08QT}{source}{}{sourcepointer-1}{}}
\newlabel{equation-three-maps@cref}{{[originalequation][0][]08QT}{[1][0][]source}{}{}{}}
\makeatother
\crefname{originallemma}{original lemma}{original lemmas}
\Crefname{originallemma}{Original lemma}{Original lemmas}
\makeatletter
\newlabel{lemma-pi-lower-shriek-polynomial-algebra}{{08QG}{source}{}{sourcepointer-2}{}}
\newlabel{lemma-pi-lower-shriek-polynomial-algebra@cref}{{[originallemma][0][]08QG}{[1][0][]source}{}{}{}}
\makeatother
\crefname{originallemma}{original lemma}{original lemmas}
\Crefname{originallemma}{Original lemma}{Original lemmas}
\makeatletter
\newlabel{simplicial-lemma-filtered-colimit-trivial-kan}{{08Q5}{source}{}{sourcepointer-3}{}}
\newlabel{simplicial-lemma-filtered-colimit-trivial-kan@cref}{{[originallemma][0][]08Q5}{[1][0][]source}{}{}{}}
\makeatother
\crefname{originallemma}{original lemma}{original lemmas}
\Crefname{originallemma}{Original lemma}{Original lemmas}
\makeatletter
\newlabel{simplicial-lemma-products-homotopy}{{08Q4}{source}{}{sourcepointer-4}{}}
\newlabel{simplicial-lemma-products-homotopy@cref}{{[originallemma][0][]08Q4}{[1][0][]source}{}{}{}}
\makeatother
\crefname{originallemma}{original lemma}{original lemmas}
\Crefname{originallemma}{Original lemma}{Original lemmas}
\makeatletter
\newlabel{simplicial-lemma-trivial-kan-base-change}{{08NN}{source}{}{sourcepointer-5}{}}
\newlabel{simplicial-lemma-trivial-kan-base-change@cref}{{[originallemma][0][]08NN}{[1][0][]source}{}{}{}}
\makeatother
\crefname{originallemma}{original lemma}{original lemmas}
\Crefname{originallemma}{Original lemma}{Original lemmas}
\makeatletter
\newlabel{simplicial-lemma-trivial-kan-homotopy}{{08NS}{source}{}{sourcepointer-6}{}}
\newlabel{simplicial-lemma-trivial-kan-homotopy@cref}{{[originallemma][0][]08NS}{[1][0][]source}{}{}{}}
\makeatother
\crefname{originallemma}{original lemma}{original lemmas}
\Crefname{originallemma}{Original lemma}{Original lemmas}
\makeatletter
\newlabel{sites-cohomology-lemma-compute-by-cosimplicial-resolution}{{08Q9}{source}{}{sourcepointer-7}{}}
\newlabel{sites-cohomology-lemma-compute-by-cosimplicial-resolution@cref}{{[originallemma][0][]08Q9}{[1][0][]source}{}{}{}}
\makeatother
\title{The lower shriek along the third projection for a polynomial algebra on top}
\author{The Stacks Project Authors}
\date{stacks-project:cotangent.tex:local-snapshot-20261012; GNU FDL 1.2 or later}
\begin{document}
\maketitle
\setcounter{subsection}{1}
\noindent\textit{Editorial note.} All delivered mathematical text is the author's own, copied from the pinned Stacks Project chapter file cotangent.tex; no statement or proof step was rewritten, reordered or removed.\par\medskip
\noindent\textit{Editorial note.} The statement refers to the three-maps equation of the chapter, which is displayed earlier in the same chapter; the proof uses only the notation it establishes.\par\medskip
\noindent\textit{Editorial note.} The seven references of the counted spans are recorded as non-load-bearing navigation pointers; the argument that consumes them stays inside the counted proof span.\par\medskip
\noindent\textit{Editorial note.} Preamble: the retained preamble of this chapter (sources/stacks-cotangent/retained-preamble.tex, lines 1-63), sliced mechanically from the chapter preamble.\par\medskip
\noindent\textit{Allowed standard prerequisites (editorial).}\begin{itemize}
\item Computation of the lower shriek via a cosimplicial resolution. Cohomology on Sites, Lemma sites-cohomology-lemma-compute-by-cosimplicial-resolution, used at proof lines 1453-1461. Reduces the claim to a homotopy statement about the mapping simplicial sets.
\item Trivial Kan fibrations: base change, homotopy equivalence and filtered colimits. Simplicial, Lemmas simplicial-lemma-trivial-kan-base-change, simplicial-lemma-trivial-kan-homotopy and simplicial-lemma-filtered-colimit-trivial-kan, used at proof lines 1482-1503. Supply the two generating cases of the homotopy computation.
\item Products of simplicial sets homotopy equivalent to a point. Simplicial, Lemma simplicial-lemma-products-homotopy, used at proof lines 1480-1481. Reduces the general case to the generating cases.
\item The value of the lower shriek on the polynomial algebra. Lemma lemma-pi-lower-shriek-polynomial-algebra, used at proof lines 1459-1461. Identifies pi\_\{3,!\}F with the constant simplicial group F(C -> C).
\end{itemize}\medskip
\ifdefined\hypertarget\hypertarget{sourcepointer-1}{}\fi
\noindent\textit{Source reference.} Original equation 08QT: focused-seeds/AGP/sources/stacks-snapshot-20261011/tags.txt, source lines 11071--11071. Navigation pointer into the Stacks Project; the standard fact it supplies is applied inside the counted proof span and the argument resting on it stays there.\par\medskip
\ifdefined\hypertarget\hypertarget{sourcepointer-2}{}\fi
\noindent\textit{Source reference.} Original lemma 08QG: focused-seeds/AGP/sources/stacks-snapshot-20261011/tags.txt, source lines 11059--11059. Navigation pointer into the Stacks Project; the standard fact it supplies is applied inside the counted proof span and the argument resting on it stays there.\par\medskip
\ifdefined\hypertarget\hypertarget{sourcepointer-3}{}\fi
\noindent\textit{Source reference.} Original lemma 08Q5: focused-seeds/AGP/sources/stacks-snapshot-20261011/tags.txt, source lines 11048--11048. Navigation pointer into the Stacks Project; the standard fact it supplies is applied inside the counted proof span and the argument resting on it stays there.\par\medskip
\ifdefined\hypertarget\hypertarget{sourcepointer-4}{}\fi
\noindent\textit{Source reference.} Original lemma 08Q4: focused-seeds/AGP/sources/stacks-snapshot-20261011/tags.txt, source lines 11047--11047. Navigation pointer into the Stacks Project; the standard fact it supplies is applied inside the counted proof span and the argument resting on it stays there.\par\medskip
\ifdefined\hypertarget\hypertarget{sourcepointer-5}{}\fi
\noindent\textit{Source reference.} Original lemma 08NN: focused-seeds/AGP/sources/stacks-snapshot-20261011/tags.txt, source lines 10996--10996. Navigation pointer into the Stacks Project; the standard fact it supplies is applied inside the counted proof span and the argument resting on it stays there.\par\medskip
\ifdefined\hypertarget\hypertarget{sourcepointer-6}{}\fi
\noindent\textit{Source reference.} Original lemma 08NS: focused-seeds/AGP/sources/stacks-snapshot-20261011/tags.txt, source lines 11000--11000. Navigation pointer into the Stacks Project; the standard fact it supplies is applied inside the counted proof span and the argument resting on it stays there.\par\medskip
\ifdefined\hypertarget\hypertarget{sourcepointer-7}{}\fi
\noindent\textit{Source reference.} Original lemma 08Q9: focused-seeds/AGP/sources/stacks-snapshot-20261011/tags.txt, source lines 11052--11052. Navigation pointer into the Stacks Project; the standard fact it supplies is applied inside the counted proof span and the argument resting on it stays there.\par\medskip
\begin{lemma}
\label{lemma-polynomial-on-top}
With notation as in (\ref{equation-three-maps})
suppose that $C$ is a polynomial algebra over $B$. Then
$L\pi_!(g_3^{-1}\mathcal{F}) = L\pi_{3, !}\mathcal{F} = \pi_{3, !}\mathcal{F}$
for any abelian sheaf $\mathcal{F}$ on $\mathcal{C}_{C/B}$
\end{lemma}

\begin{proof}
Write $C = B[E]$ for some set $E$. Choose a resolution
$P_\bullet \to B$ of $B$ over $A$. For every $n$ consider
the object $U_n = (P_n \to B, P_n[E] \to C)$ of $\mathcal{C}_{C/B/A}$.
Then $U_\bullet$ is a cosimplicial object of $\mathcal{C}_{C/B/A}$.
Note that $u_3(U_\bullet)$ is the constant cosimplicial
object of $\mathcal{C}_{C/B}$ with value $(C \to C)$.
We will prove that the object $U_\bullet$ of $\mathcal{C}_{C/B/A}$
satisfies the hypotheses of
Cohomology on Sites, Lemma
\ref{sites-cohomology-lemma-compute-by-cosimplicial-resolution}.
This implies the lemma as it shows that $L\pi_!(g_3^{-1}\mathcal{F})$
is computed by the constant simplicial abelian group
$\mathcal{F}(C \to C)$ which is the value of
$L\pi_{3, !}\mathcal{F} = \pi_{3, !}\mathcal{F}$ by
Lemma \ref{lemma-pi-lower-shriek-polynomial-algebra}.

\medskip\noindent
Let $U = (\beta : P \to B, \gamma : Q \to C)$ be an object of
$\mathcal{C}_{C/B/A}$. We may write $P = A[S]$ and $Q = A[S \amalg T]$
by the definition of our category $\mathcal{C}_{C/B/A}$. We have to show that
$$
\Mor_{\mathcal{C}_{C/B/A}}(U_\bullet, U)
$$
is homotopy equivalent to a singleton simplicial set $*$. Observe that this
simplicial set is the product
$$
\prod\nolimits_{s \in S} F_s \times \prod\nolimits_{t \in T} F'_t
$$
where $F_s$ is the corresponding simplicial set for
$U_s = (A[\{s\}] \to B, A[\{s\}] \to C)$
and $F'_t$ is the corresponding simplicial set for
$U_t = (A \to B, A[\{t\}] \to C)$. Namely, the object $U$
is the product $\prod U_s \times \prod U_t$ in $\mathcal{C}_{C/B/A}$.
It suffices each $F_s$ and $F'_t$ is homotopy equivalent to $*$, see
Simplicial, Lemma \ref{simplicial-lemma-products-homotopy}.
The case of $F_s$ follows as $P_\bullet \to B$ is a trivial Kan
fibration (as a resolution) and $F_s$ is the fibre of this map over
$\beta(s)$. (Use Simplicial, Lemmas
\ref{simplicial-lemma-trivial-kan-base-change} and
\ref{simplicial-lemma-trivial-kan-homotopy}).
The case of $F'_t$ is more interesting. Here we are saying that
the fibre of
$$
P_\bullet[E] \longrightarrow C = B[E]
$$
over $\gamma(t) \in C$ is homotopy equivalent to a point. In fact we
will show this map is a trivial Kan fibration. Namely,
$P_\bullet \to B$ is a trivial can fibration. For any ring $R$
we have
$$
R[E] =
\colim_{\Sigma \subset \text{Map}(E, \mathbf{Z}_{\geq 0})\text{ finite}}
\prod\nolimits_{I \in \Sigma} R
$$
(filtered colimit). Thus the displayed map of simplicial sets is a
filtered colimit of trivial Kan fibrations, whence a trivial Kan fibration
by Simplicial, Lemma \ref{simplicial-lemma-filtered-colimit-trivial-kan}.
\end{proof}

\end{document}
